%==============================================================
%  Curriculum vitae — Aleksandra (Sasha) Erakhtina
%
%  Compile:  pdflatex cv.tex     (run twice, for page numbers)
%  Typeface: Computer Modern (the LaTeX default — no packages needed)
%
%  \phonetrue  prints the mobile number — the copy you email people
%  \phonefalse omits it              — the copy you put on the website
%
%  Items still to confirm are marked  \todo{...}  and print in the
%  margin so you cannot miss them. Delete each one as you fill it,
%  then delete the \todo macro itself before circulating.
%==============================================================

\documentclass[11pt,a4paper]{article}

\newif\ifphone
\phonefalse

\newif\ifreferees
\refereesfalse       % \refereestrue prints referees with their email addresses

\usepackage[a4paper,top=2.2cm,bottom=2.2cm,left=2.4cm,right=2.4cm]{geometry}
\usepackage[expansion=false]{microtype}
\usepackage{titlesec}
\usepackage{enumitem}
\usepackage{fancyhdr}
\usepackage{needspace}
\usepackage[hidelinks]{hyperref}



\setlength{\parindent}{0pt}

%---- placeholder marker (delete this macro before circulating) ----
\newcommand{\todo}[1]{\textbf{[#1]}}

%---- your GitHub username: set it once, used in the header ----
\newcommand{\username}{sasha-erakhtina}

%---------------- section headings ------------------------------
\titleformat{\section}
  {\normalfont\normalsize\scshape}
  {}{0pt}{}[\vspace{1pt}\hrule height 0.4pt]
\titlespacing*{\section}{0pt}{16pt}{7pt}

%---------------- dated entries ---------------------------------
\newcommand{\entry}[2]{%
  \begin{minipage}[t]{0.15\textwidth}\raggedright\small #1\end{minipage}%
  \hfill%
  \begin{minipage}[t]{0.83\textwidth}#2\end{minipage}%
  \par\vspace{5pt}}

% an entry with no date column — used for the working papers
\newcommand{\nodate}[1]{#1\par\vspace{7pt}}

% one-line squeeze: slightly smaller and tighter, so a long title stays on one line
\newcommand{\oneline}[1]{{\footnotesize\raggedright #1\par}}

% \paper{title}{coauthors}{description}
\newcommand{\paper}[3]{%
  \emph{``#1''}\par\vspace{1pt}
  \ifx\relax#2\relax\else{\small #2}\par\vspace{1pt}\fi
  {\small #3}\par}

% role sub-heading inside a section (Lecturer / Teaching assistant)
\newcommand{\subsec}[1]{\vspace{2pt}{\scshape\small #1}\par\vspace{3pt}}

\newlist{tight}{itemize}{1}
\setlist[tight]{leftmargin=1.2em, itemsep=2pt, topsep=2pt, parsep=0pt, label=--}

\pagestyle{fancy}
\fancyhf{}
\renewcommand{\headrulewidth}{0pt}
% \fancyfoot[L]{\small last updated: 08/2026}
\fancyfoot[R]{\small \thepage}

%================================================================
\begin{document}

%---------------- masthead --------------------------------------

{\large \textbf{Aleksandra (Sasha) Erakhtina}}\\[5pt]
  { \textbf{Email:} \href{mailto:aleksandra.erakhtina@uts.edu.au}{aleksandra.erakhtina@uts.edu.au}}\\
  { \textbf{Website:} \href{https://\username.github.io}{\username.github.io}}%
  \ifphone\\
  {\small \textbf{Mobile:} +61 451 954 050}\fi\\
  { \textbf{Address:} 14-28 Ultimo Rd, Ultimo NSW 2007
Australia
}
\hyphenpenalty=10000
\exhyphenpenalty=10000
\tolerance=2000
\emergencystretch=3em


%---------------- positions -------------------------------------
\section{Academic Positions}

\entry{2026--present}{\textbf{Postdoctoral Research Associate}\par
  Economics Discipline Group, UTS Business School, Australia\par
  {\small \textit{ARC Discovery Project on forced migration and its effects across generations}}}


%---------------- education -------------------------------------
\section{Education}

\entry{2021--2026}{\textbf{PhD in Economics}, 
  University of Technology Sydney, Australia\par
  \oneline{Thesis: \textit{Bound by Constraints, Guided by Expectations: Essays in Economic Decision-Making}}\par
  { Principal supervisor: Adeline Delavande}\par
  { Co-supervisors: Elif Incekara-Hafalir and Esther Mirjam Girsberger}}

\entry{2017--2020}{\textbf{Postgraduate Degree in Economics}\par
  Institute of Economics and Industrial Engineering, Russia}

\entry{2016--2017}{\textbf{Master of Economics}, Analyse et Politique Économique (M2) \par
  Paris School of Economics, France}

\entry{2015--2017}{\textbf{Master of Economics}, Quantitative Economics\par
  Novosibirsk State University, Russia}

\entry{2010--2015}{\textbf{Bachelor of Management}, Honours\par
  Novosibirsk State University, Russia}


%---------------- working papers --------------------------------
\section{Working Papers}

\nodate{\textbf{\paper
  {Sex Ratio, Marriage Market and Assortative Matching in Colonial Tasmania}}
  {\relax}
  {Convict transportation left nineteenth-century Tasmania with a heavily male-biased sex ratio for decades. This setting allows a rare look at how such an imbalance shaped the marriage market, both while it persisted and long after. Combining individual marriage records with historical census data, we use the administrative allocation of convicts, conditional on labour demand, as a quasi-random source of variation in local sex ratios. The results reveal a two-stage response. In the short term, during the transportation era, male-biased sex ratios directly influenced match formation: older women were more likely to marry, and assortative matching on literacy and convict status weakened. In the longer term, over the fifty years following transportation, these distortions settled into both marriage and family formation. Marital-status sorting shifted in women's favour, and widows were more likely to remarry in the early post-transportation years. Couples in historically more male-biased districts went on to have substantially larger families, and this effect, unlike the others, strengthened rather than faded over successive cohorts, persisting well after the imbalance itself had eased. Underlying mechanisms plausibly include better spouse quality, lower female labour supply, and more stable unions. A complementary analysis of convict registers shows that marriage under such imbalance was highly selective: men’s marriage chances depended on youth, literacy, stature, and occupational skills, whereas women’s prospects were shaped more by personal attributes such as appearance, origin, and offending history.}}
\newpage
\nodate{\textbf{\paper
  {The Long-Term Effects of Historical Gender Imbalances on Children's Educational Achievement in Australia}}
  {\relax}
  {Can a demographic shock leave a mark on human capital more than a century after it has faded? This paper examines the long-run effect of a historically male-biased sex ratio on children's educational achievement in Australia. The analysis exploits quasi-random variation in local sex ratios arising from the nineteenth-century convict transportation system, which allocated overwhelmingly male convicts across districts according to labour demand, creating large and uneven gender imbalances. By linking historical census data to contemporary longitudinal records on children, the analysis uses the convict sex ratio as an instrument for the total population historical sex ratio. The results show that children living in areas with more male-biased populations in the 1840s achieve significantly lower test scores in numeracy and writing today - by roughly 9 \% of a standard deviation at the end of compulsory schooling. These effects are concentrated among boys, echoing the gendered nature of the original shock, and persist after accounting for detailed individual, family, school, and neighbourhood controls. Evidence on mechanisms is consistent with weaker educational incentives for men and enduring local norms surrounding gender and masculinity in historically male-biased areas.}}

\nodate{\paper
  {\textbf{When Do You Expect to Be Paid? Expectations Shape Patience}}
  {\textit{with Elif Incekara-Hafalir and Benjamin Young}}
   {\relax}
   {Time preferences shape decisions about saving, education, and health, yet research has largely overlooked how expectations about payment timing influence elicited measures of time preferences. We provide causal evidence that such expectations systematically affect intertemporal choices. In a pre-registered experiment, we manipulate payment-timing expectations while holding choice menus fixed. Participants expecting payment in the future choose more patient options than those expecting payment today, despite identical rewards and delays. A reference-dependent model, where expected timing serves as a temporal reference point, explains these shifts. The findings illuminate patterns in the time preference literature and suggest a low-cost policy lever.}}


\section{Work in Progress}

\nodate{\paper
  {\textbf{Impact of Forced Migration on Tasmanian Convicts and Their Descendants}}
  {\textit{with Adeline Delavande, Luis Pontes de Vasconcelos,
   Hamish Maxwell-Stewart and Ran Abramitzky}}
      {\relax}
  {This project evaluates the impact of transportation on the life outcomes of convicts and their
   descendants by comparing transported prisoners with named brothers and sisters who remained in
   Britain and Ireland, so that unobserved family background and shared traits are held fixed. It
   then examines how punishment and place shaped outcomes within Tasmania, and traces how these effects carried across generations. The project extends existing digitisations and
   record linkages for the colonial period, and offers a longitudinal perspective on place-based
   conditions relevant to contemporary migration and settlement policy.}}

\nodate{\paper
  {\textbf{Expectations and Time Preferences: A Within-Subject Design}}
  {\textit{with Elif Incekara-Hafalir, Benjamin Young and Xueting Wang}}
     {\relax}
  {Erakhtina et al (2026) provided causal evidence that expectations systematically affect intertemporal choices. Specifically, participants who expected payment in the future chose more patient options than participants who expected payment today, despite facing identical amounts, delays, and choice menus. We report a within-subject test of that mechanism. The payment is expected to be paid today in one part and in eight weeks in the other, with order counterbalanced. Participants make choices between money today and money in eight weeks, plus lottery choices that test whether the reference date shifts risk attitudes as well. Because the same individual is observed under both reference dates, the design allows us to examine within-person changes in choice as expected reward timing changes. We will present evidence on whether individuals’ patience changes across the two expected payment timings, as well as heterogeneity in these responses to identify who is more likely to change their intertemporal choices when expectations about reward timing differ.
}}

\nodate{\paper
  {\textbf{Improving the Utility of Private Health Insurance Choice in Dual Public--Private Markets}}
  {\textit{with Nathan Kettlewell and Hong Il Yoo}}
{\relax}
     {We study the quality of health insurance choices when private insurance offers a quality advantage (shorter waiting times) rather than reduced financial risk. In a controlled experiment modelled on the Australian system, we elicit willingness to pay to avoid waiting, risk preferences from incentivised lotteries, and insurance plan choices with and without decision aids. Choice quality is defined relative to welfare-maximising benchmarks that integrate monetary, risk, and waiting-time preferences. Our analysis combines two approaches. First, we estimate preferences separately for lottery and insurance tasks, allowing correlation across them. Second, we impose a coherent structural model in which utility curvature is common across tasks and waiting-time attributes act as source-dependent shifts in probability weighting. We then evaluate decision aids using both non-parametric tests (dominated choices) and structural welfare measures (consumer surplus losses relative to preference-consistent benchmarks). We find that participants often fail to incorporate waiting-time preferences into insurance choices, leading to systematic welfare losses. Decision aids improve outcomes, with effectiveness varying by tool. The results demonstrate the cognitive difficulty of integrating monetary and non-monetary attributes in insurance decisions and provide guidance for designing interventions that improve consumer choice quality.}}
     
%---------------- teaching --------------------------------------
\section{Teaching}

\entry{2023--2026}{University of Technology Sydney (teaching assistant)\par
  {\small 26134 Business Statistics; 23941 Mathematics for Economists;
   23508 Quantitative Methods in Economics and Business.}}

\entry{2018--2022}{Novosibirsk State University (teaching assistant, lecturer)\par
  {\small Experimental Economics;
   Experimental Economics: Programming in oTree;
   Data Analysis in R.}}



   %---------------- research assistance ---------------------------
\section{Research Assistance}

\entry{2022--2026}{\textbf{University of Technology Sydney}\par
  {\small Adeline Delavande, \emph{The Ambivalence about Pregnancy}, \emph{Covid Ambiguity},
   \emph{Uncertainty Attitudes}; with Peter Siminski, \emph{Impact Evaluation} (Paul Ramsay
   Foundation). Elif Incekara-Hafalir, \emph{A Family of Norms: Related Beliefs about Female
   Labor Force Participation}. Anne Summers, \emph{Domestic Violence and Employment}, and more.}}

\entry{2017--2020}{\textbf{Institute of Economics and Industrial Engineering}, Russia}

%---------------- presentations ---------------------------------
\section{Conference Presentations and Seminars}
\newcommand{\pres}[2]{#1\par{\small\emph{#2}}\par\vspace{5pt}}
\entry{2026}{%
  \pres{Professor Lucy Frost: Digitising the Colonial Past, Annual Symposium}
       {Sex Ratio, Marriage Market and Assortative Matching. Evidence from Colonial Tasmania}
  \pres{SIOE, Society for Institutional \& Organizational Economics Conference, Fontainebleau}
       {The long-term effects of historical gender imbalances on children's educational achievement in Australia}
  \pres{Asia-Pacific Economic and Business History Meeting, Adelaide}
       {The long-term effects of historical gender imbalances on children's educational achievement in Australia}}

\entry{2025}{%
  \pres{Wellington Applied Econometrics Workshop, New Zealand}
       {Marriage Market and Assortative Matching in Colonial Tasmania}
  \pres{SIOE, Society for Institutional \& Organizational Economics Conference, Sydney}
       {Marriage Market and Assortative Matching in Colonial Tasmania}
  \pres{Death, Demography and Digital History Symposium, Melbourne}
       {Marriage Market and Assortative Matching in Colonial Tasmania}
  \pres{The Economic History Society Annual Conference, Glasgow}
       {The long-term effects of historical gender imbalances on adolescents in Australia}
  \pres{AGEW, Australian Gender Economics Workshop, Wollongong}
       {Marriage Market and Assortative Matching in Colonial Tasmania}}

\entry{2024}{%
  \pres{ANEW, 11th Annual Natural Experiments Workshop, Melbourne}
       {Marriage Market and Assortative Matching in Colonial Tasmania}
  \pres{Microdata in Economic History: Beyond the Full-Count Census, Munich}
       {Marriage Market and Assortative Matching in Colonial Tasmania}
  \pres{UTS Behavioural Lab Meeting, Sydney}
       {Private health insurance choice in dual public-private markets}
  \pres{UTS Behavioural Lab Conference, Sydney}
       {Colonisation, convicts and marriage market in Tasmania}}

\entry{2023}{%
  \pres{UTS Behavioural Lab Meeting, Sydney}
       {The Impact of Expectations on Time Preferences}
  \pres{UTS Behavioural Lab Conference, Sydney}
       {The Impact of Expectations on Time Preferences}}
%---------------- awards ----------------------------------------
\section{Awards, Scholarships and Funding}


\entry{2021--2026}{Ross Milbourne Research Scholarship; International Research Scholarship, UTS}
\entry{2025}{UTS Vice-Chancellor's Conference Fund}

\entry{2020}{Award for the development of the \emph{Behavioural Economics} Master's program, Novosibirsk State University}
  
\entry{2018}{Best paper presentation, XIV Autumn Conference of Young Scientists (Russia)}



\entry{2016, 2017}{Best paper presentation, ISSC Interdisciplinary Studies (Russia)}

\entry{2016--2017}{French Government Scholarship}

\entry{2014--2015}{Baker Hughes Scholarship}

%---------------- service ---------------------------------------
\section{Academic Service and Conference Organisation}

\entry{2025}{Programme support, European Society for Population Economics Annual Conference;
  programme support, Society for Institutional and Organizational Economics Annual Conference.}

\entry{2024}{Discussant, Australian Cliometrics Workshop, UNSW; organising support, UTS Behavioural
  Lab Conference.}

\entry{2023}{Organising team, Australian Gender Economics Workshop; organising support, UTS
  Behavioural Lab Conference.}

\entry{2018--2020}{Organising \& programme support, Conference of Young Scientists, Russia.}

%---------------- workshops -------------------------------------
\section{Workshops and Training}

\entry{2025}{Training Course on Coding Historical Causes of Death}

\entry{2020--2024}{ISEO Summer School, \emph{The World After: Challenges Ahead for the Global
  Economy}, Italy; 24th Visiting Graduate Student Workshop in Experimental Economics, ESI Chapman
  University, USA; Summer School on Socioeconomic Inequality, HCEO University of Chicago and
  New Economic School; Summer School on Experimental Economics, Higher School of Economics, Russia;
  XII Russian Summer School on Institutional Analysis, Higher School of Economics, Russia;
  V Winter Scientific School on Diversity, New Economic School, Russia.}



%---------------- software --------------------------------------
\section{Software and languages}

\entry{Programming}{R, Python, Stata, Matlab}

\entry{Spatial}{QGIS, ArcGIS (with R) - georeferencing and digitising historical maps}

\entry{Experiments}{oTree \& Qualtrics, Prolific \& Pureprofile}

\entry{Web}{HTML/CSS, JavaScript (integration with Qualtrics)}





\entry{Languages}{Russian, English (fluent), German (intermediate), French (beginner)}
%---------------- referees --------------------------------------
\ifreferees
\needspace{4\baselineskip}
\section{Referees}

\nodate{\textbf{Prof.\ Adeline Delavande} \par
  {\small University of Technology Sydney\ \ $\cdot$\ \
   \href{mailto:Adeline.Delavande@uts.edu.au}{Adeline.Delavande@uts.edu.au}}}

\nodate{\textbf{Prof. Elif Incekara-Hafalir} \par
  {\small University of Technology Sydney\ \ $\cdot$\ \
   \href{mailto:Elif.IncekaraHafalir@uts.edu.au}{Elif.IncekaraHafalir@uts.edu.au}}}

\nodate{\textbf{Dr Esther Mirjam Girsberger} \par
  {\small University of Technology Sydney\ \ $\cdot$\ \
   \href{mailto:EstherMirjam.Girsberger@uts.edu.au}{EstherMirjam.Girsberger@uts.edu.au}}}
\fi

\end{document}
